%=============================================================================!
% 2.5  WEIGHT, FREE FALL AND FREE-BODY DIAGRAMS                               !
%=============================================================================!
\section{Weight, free fall and free-body diagrams}
\label{sec:ne-gravity}

With the second law in hand we can explain the first motion of
Chapter~1, the thrown ball, which was there described without a cause.

\subsection*{Weight}

Near the surface of the Earth every body is pulled downwards by a force,
its weight. Hanging a body at rest from a spring balance measures its
weight: the body does not accelerate, so by the second law the balance's
upward pull equals the Earth's downward pull (the lamp below works this
through). Such measurements show that the weight is proportional to the
mass: twice, three times or half the mass gives twice, three times or
half the reading. We therefore write the weight as $m\vb{g}$, where $\vb{g}$ is a vector
pointing straight down whose length $g$ is the same for every body. Since
weight divided by mass is in newtons per kilogram, which by the definition
of the newton is \si{\metre\per\second\squared}, $g$ has the units of an
acceleration. Its standard value is $g = \SI{9.80665}{\metre\per\second
\squared}$, and the free-fall argument below shows that it is the $g$ of
Chapter~1. An apple of \SI{100}{\gram} weighs $0.1\times
9.80665 = \SI{0.981}{\newton}$, about one newton.

In everyday speech weight often means mass, as in `a bag of flour weighs
a kilogram'. This book keeps them apart. Mass, in kilograms, measures
inertia and is the same everywhere; weight, in newtons, is the pull of the
Earth, and on the Moon, whose pull is weaker, the same bag would weigh
about a sixth as much. A kitchen scale measures the force and divides by
$g$ to display a mass.

\subsection*{Free fall}

A ball in flight, with air resistance negligible, feels its weight and
nothing else, a motion called free fall. The second law gives $m\vb{a} =
m\vb{g}$, and dividing both sides by $m$,
\begin{equation*}
   \vb{a} = \vb{g} .
\end{equation*}
The mass cancels, so every body falls with the same acceleration, whatever
its mass. This is why $g$ in Chapter~1 was the same for every object, and
it rests on the experimental fact that weight is proportional to mass. A
feather and a coin dropped together in a room do not land together, but
only because the drag of the air (\cref{sec:ne-drag}) is large compared
with the feather's tiny weight. Without air they fall together, as an
astronaut showed on the Moon in 1971 by dropping a hammer and a feather
side by side.

We can now derive the motion of the ball that Chapter~1 took from
observation. The motion is along a single vertical line, so, as in
\cref{eq:mo-ball}, we call the height $x$ and measure it upwards. The
component of $\vb{a} = \vb{g}$ along this line is
\begin{equation*}
   \frac{\dd^2 x}{\dd t^2} = -g ,
\end{equation*}
the minus sign because $\vb{g}$ points down. Integrating once from $t =
0$, as in \cref{sec:mo-integrating}, gives $\dd x/\dd t = v_0 - gt$, where
$v_0$ is the launch velocity. Integrating again, with the ball leaving the
hand at $x = 0$, gives
\begin{equation*}
   x(t) = v_0 t - \tfrac12 g t^2 ,
\end{equation*}
which is \cref{eq:mo-ball}. The projectile follows in the same way: the
horizontal component of $\vb{g}$ is zero, so the horizontal velocity
stays constant.

For atoms, gravity hardly matters: \cref{sec:ne-atoms} shows that the
weight of an atom is some $10^{16}$ times smaller than the forces its
neighbours exert on it.

\subsection*{Free-body diagrams}

Applying the second law to a body starts with a free-body diagram: a
sketch of the body alone, with an arrow for every force acting on it and
for nothing else. \Cref{fig:ne-free-body}(a) shows the ball in flight,
which feels only its weight; its velocity is drawn dashed, because a
velocity is not a force. \Cref{fig:ne-free-body}(b) shows a lamp of
\SI{2}{\kilogram} hanging at rest from a cord. Two forces act on it: its
weight, $mg$ downwards, and the pull of the cord, called its tension,
upwards, which we write $F_{\mathrm{cord}}$. The lamp does not accelerate,
so by the second law the net force on it is zero, and taking upwards as
positive,
\begin{equation*}
   F_{\mathrm{cord}} - mg = 0
   \quad\Longrightarrow\quad
   F_{\mathrm{cord}} = mg = 2\times9.80665\ \si{\newton}
   = \SI{19.6}{\newton} .
\end{equation*}
The arrows of a free-body diagram are labelled with the sizes of the
forces, and their directions are shown by the arrows themselves.

\begin{figure}[htb]
\centering
\begin{tikzpicture}[scale=1.25, line cap=round,
   force/.style={-{Stealth[length=6pt]}, line width=1.2pt},
   body/.style={line width=0.6pt, fill=black!8}]
   % (a) the ball in flight
   \begin{scope}
      \node at (-1.3,1.55) {\small (a)};
      \draw[body] (0,0) circle (0.3);
      \draw[accent, -{Stealth[length=5pt]}, dashed, line width=0.8pt]
         (0.3,0.17) -- (1.0,0.85) node[right] {\small $\vb{v}$};
      \draw[force] (0,0) -- (0,-1.25) node[right] {\small $mg$};
      \fill (0,0) circle (1.4pt);
      \node at (0,-1.75) {\small ball in flight};
   \end{scope}
   % (b) the lamp hanging at rest from a cord
   \begin{scope}[xshift=4.0cm]
      \node at (-1.3,1.55) {\small (b)};
      \fill[black!35] (-0.6,1.45) rectangle (0.6,1.55);
      \draw[line width=0.8pt] (0,1.45) -- (0,0.3);
      \draw[body] (-0.35,0.3) -- (0.35,0.3) -- (0.5,-0.2) -- (-0.5,-0.2)
         -- cycle;
      \draw[force] (0,0.05) -- (0,1.2)
         node[midway, right=2pt] {\small $F_{\mathrm{cord}}$};
      \draw[force] (0,0.05) -- (0,-1.25) node[right] {\small $mg$};
      \fill (0,0.05) circle (1.4pt);
      \node at (0,-1.75) {\small hanging lamp};
   \end{scope}
\end{tikzpicture}
\caption{Free-body diagrams. (a) A ball in flight feels only its weight;
its velocity is dashed, because it is not a force. (b) A lamp hanging at
rest feels its weight and the tension $F_{\mathrm{cord}}$ of the cord,
which balance.}
\label{fig:ne-free-body}
\end{figure}

\begin{selfcheck}
A hammer and a feather are dropped together, first in a vacuum, then in
air. Which lands first in each case, and why?
\end{selfcheck}
